%====================================
% KAPITEL 4: GLIEDERUNG UND ZEITPLANUNG
%====================================
\section{Vorl\"aufige Gliederung der Bachelor-Thesis}

Die Bachelorarbeit gliedert sich in sechs Hauptkapitel, die den Forschungsprozess von der Problemstellung bis zur abschlie\ss{}enden Bewertung nachvollziehbar abbilden. Bei einem Gesamtumfang von rund 38 bis 45 Seiten bleibt dabei ausreichend Raum f\"ur eine fundierte theoretische Fundierung, eine transparente Darstellung des methodischen Vorgehens sowie eine differenzierte Diskussion der Ergebnisse.

Den Auftakt bildet die \textbf{Einleitung}, die in die Problemstellung einf\"uhrt, die Relevanz des Themas begr\"undet und die zentralen Forschungsfragen ableitet. Im daran anschlie\ss{}enden Kapitel der \textbf{Theoretischen Grundlagen} werden die relevanten Fachkonzepte und der aktuelle Forschungsstand auf Basis einer systematischen Literaturrecherche aufbereitet. Das dritte Kapitel \textbf{Methodisches Vorgehen} legt offen, nach welchen Kriterien die Literatur recherchiert, ausgew\"ahlt und analysiert wurde, und schafft damit die notwendige Nachvollziehbarkeit f\"ur die folgende Untersuchung.

Auf dieser Grundlage pr\"asentiert das vierte Kapitel die \textbf{Ergebnisse} der Analyse, die im f\"unften Kapitel zu konkreten \textbf{Handlungsempfehlungen} verdichtet werden. Den Abschluss der Arbeit bildet das \textbf{Fazit}, das die zentralen Erkenntnisse zusammenfasst, Limitationen benennt und einen Ausblick auf weiterf\"uhrende Forschungsfragen gibt. Die nachfolgende Tabelle zeigt die vorl\"aufige Gliederung im \"Uberblick.

\begin{table}[H]
\centering
\begin{tabularx}{\textwidth}{lXr}
\toprule
\textbf{Kapitel / Abschnitt} & & \textbf{ca. Seiten} \\
\midrule
Verzeichnisse (Inhalt, Abk\"urzungen, Abb., Tab.) & & i--v \\
\midrule
1. Einleitung & & 3--4 S. \\
\quad 1.1 Problemstellung und Relevanz & & \\
\quad 1.2 Zielsetzung und Forschungsfragen & & \\
\quad 1.3 Aufbau der Arbeit & & \\
\midrule
2. Theoretische Grundlagen & & 12--13 S. \\
\quad 2.1 Enterprise-Resource-Planning-Systeme & & \\
\quad 2.2 Erfolg von ERP-Implementierungen (DeLone \& McLean) & & \\
\quad 2.3 Change Management (ADKAR-Modell) & & \\
\quad 2.4 Vorl\"aufiger Bezugsrahmen der Untersuchung & & \\
\midrule
3. Methodisches Vorgehen & & 5--6 S. \\
\quad 3.1 Systematische Literaturanalyse (Webster \& Watson) & & \\
\quad 3.2 Suchstrategie, Datenbanken und Schlagworte & & \\
\quad 3.3 Auswahlprozess und Auswertungslogik (Mayring) & & \\
\midrule
4. Ergebnisse der Literaturanalyse & & 10--12 S. \\
\quad 4.1 Deskriptive Auswertung des Samples & & \\
\quad 4.2 Identifizierte CM-Faktoren (Kommunikation, F\"uhrung, Schulung) & & \\
\quad 4.3 Verkn\"upfung mit DeLone \& McLean-Dimensionen & & \\
\midrule
5. Handlungsempfehlungen f\"ur den Mittelstand & & 5--6 S. \\
\quad 5.1 ADKAR-basierte Empfehlungen & & \\
\quad 5.2 Kritische W\"urdigung & & \\
\midrule
6. Fazit und Ausblick & & 3--4 S. \\
\quad 6.1 Beantwortung der Forschungsfragen & & \\
\quad 6.2 Limitationen und Implikationen & & \\
\midrule
Literaturverzeichnis & & \\
Anhang (PRISMA-Flowchart, Studienliste, Kodierleitfaden) & & \\
\midrule
\textbf{Gesamtumfang Textteil} & & \textbf{38--45 S.} \\
\bottomrule
\end{tabularx}
\caption{Vorl\"aufige Gliederung der Bachelor-Thesis}
\label{tab:gliederung}
\end{table}

\section{Zeitplanung der Thesis}

Die Arbeit ist auf einen Bearbeitungszeitraum von zw\"olf Wochen ausgelegt, der in drei aufeinander aufbauende Phasen gegliedert ist. Diese Strukturierung gew\"ahrleistet, dass sowohl konzeptionelle Vorarbeiten als auch die analytische Hauptarbeit sowie die abschlie\ss{}ende Verschriftlichung ausreichend Ber\"ucksichtigung finden. Der Zeitplan ist bewusst realistisch gehalten und enth\"alt an mehreren Stellen zeitliche Puffer, um auf unvorhergesehene Verz\"ogerungen angemessen reagieren zu k\"onnen.

Die erste Phase (Woche 1 bis 2) umfasst die formelle Anmeldung der Arbeit, die Fertigstellung des Expos\'es sowie die finale Festlegung der methodischen Vorgehensweise. In dieser Phase werden die theoretischen Grundlagen gesch\"arft und der Untersuchungsrahmen endg\"ultig abgesteckt. Die zweite und umfangreichste Phase (Woche 3 bis 8) bildet die Hauptarbeitsphase: In sechs Wochen erfolgen die systematische Literaturrecherche und -auswertung, die Quellenanalyse sowie die anschlie\ss{}ende Synthese der gewonnenen Erkenntnisse. Diese Phase ist bewusst am l\"angsten bemessen, da sie das analytische Kernst\"uck der Arbeit darstellt. Die dritte Phase (Woche 9 bis 12) ist der Verschriftlichung, der Feinkorrektur und der finalen Abgabe gewidmet.

Die nachfolgende Grafik visualisiert den zeitlichen Ablauf als Gantt-Diagramm.

\begin{figure}[H]
\centering
\begin{ganttchart}[y unit title=0.5cm,
y unit chart=0.6cm,
vgrid,hgrid,
title label font=\footnotesize,
title height=1,
bar/.style={fill=gray!40},
bar height=0.6,
bar label font=\footnotesize,
group/.style={draw=black, fill=black!20},
group height=0.4,
group label font=\footnotesize\bfseries
]{1}{12}
\gantttitle{Woche}{12} \\
\gantttitlelist{1,...,12}{1} \\
\ganttgroup{Phase 1: Vorbereitung}{1}{2} \\
\ganttbar{Anmeldung \& Expos\'e}{1}{2} \\
\ganttbar{Methodik-Festlegung}{2}{2} \\
\ganttgroup{Phase 2: Hauptarbeitsphase}{3}{8} \\
\ganttbar{Literaturrecherche}{3}{5} \\
\ganttbar{Auswahl \& Kodierung}{5}{6} \\
\ganttbar{Analyse \& Synthese}{7}{8} \\
\ganttgroup{Phase 3: Schreibphase}{9}{12} \\
\ganttbar{Grundlagen \& Methodik}{9}{10} \\
\ganttbar{Ergebnisse \& Empfehlungen}{10}{11} \\
\ganttbar{Korrektur \& Lektorat}{11}{12} \\
\ganttbar{Abgabe}{12}{12} \\
\end{ganttchart}
\caption{Gantt-Diagramm der Zeitplanung}
\label{fig:zeitplan}
\end{figure}
